\section{Independent-cap compactness and the complete model spectrum}
\label{sec:model-compactness}

The estimates in this section concern forms of degree one or two.
The distinction matters at the middle degree: we do not assert the
same weighted estimate for degree three or for arbitrary mixed-degree
forms.

\subsection{A weighted estimate on one completed cap}

Let $Z$ be a completed cap diffeomorphic to $T^*S^3$, with fixed smooth
complete metric equal to $\dd r^2+r^2h$ for $r\ge r_0$. Choose a
positive smooth radius $\rho_Z$ equal to $r$ on the end and bounded
above and below on the core. Put $\mathcal D_Z=\dd+\delta$.

\begin{proposition}\label{prop:weighted-cap}
For $p=1,2$, every compactly supported smooth $p$-form $v$ satisfies
\begin{equation}\label{eq:weighted-cap}
 \|v\|_{L^2(Z)}\le C_p\|\rho_Z\mathcal D_Zv\|_{L^2(Z)}.
\end{equation}
The constant depends only on the fixed cap.
\end{proposition}
\begin{proof}
On the cone the radial unitary transformation gives
\begin{equation}\label{eq:radial-gauss-bonnet}
 U\mathcal D U^{-1}
       =J\left(\partial_r+\frac{\mathsf A}{r}\right),
 \qquad
 \|u\|^2=\int\|\sigma(r)\|_{\mathcal H}^2\,\dd r,
 \quad \sigma=Uu.
\end{equation}
Here $\mathcal H$ is the doubled space of all link forms, $J$ is
unitary, and $\mathsf A$ is self-adjoint elliptic with compact
resolvent. We use the normalization in
\cite[Section~2.2, p.~1028]{anne-takahashi}. For a pure $p$-form,
the radial and tangential entries belong to the closed subspace
$\mathcal S_p=L^2\Omega^{p-1}(L)\oplus L^2\Omega^p(L)$.
It need not be invariant under $\mathsf A$.

Set $t=\log r$, $w(t)=r^{1/2}\sigma(r)$, and
$\mathsf B=\mathsf A-\frac12$. Then
\begin{equation}\label{eq:log-radial-norms}
 \|u\|^2=\int\|w(t)\|^2\,\dd t,\qquad
 \|r\mathcal Du\|^2
       =\int\|(\partial_t+\mathsf B)w(t)\|^2\,\dd t.
\end{equation}
We claim that
\begin{equation}\label{eq:pure-critical-kernel}
 \mathcal S_p\cap\ker\mathsf B=\{0\}\qquad(p\ne3).
\end{equation}
A vector in this intersection would give a pure harmonic form
$u=r^{p-3}\alpha+\dd r\wedge r^{p-4}\beta$, of pointwise order
$r^{-3}$. Because it has a single degree, $\dd u$ and
$\delta u$ have distinct degrees and vanish separately. Their
equations are
\[
 \dd_L\alpha=0,\quad (p-3)\alpha=\dd_L\beta,\qquad
 \delta_L\beta=0,\quad \delta_L\alpha+(p-3)\beta=0.
\]
It follows that $\Delta_L\beta=-(p-3)^2\beta$. Nonnegativity
forces $\beta=0$ and then $\alpha=0$ for $p\ne3$. Kernel vectors
are smooth by ellipticity, so this proves the claim on the full
kernel.

Since $\mathsf B$ has a spectral gap off its finite-dimensional
kernel, \eqref{eq:pure-critical-kernel} implies
\begin{equation}\label{eq:pure-tangential-coercivity}
 \|b\|\le C_p\|\mathsf Bb\|,\qquad
 b\in\mathcal S_p\cap\Dom\mathsf B.
\end{equation}
Indeed, if unit vectors in this subspace had $\mathsf Bb_j\to0$,
their nonkernel components would tend to zero. Their kernel
components have a convergent subsequence, giving a nonzero limit
in both $\mathcal S_p$ and $\ker\mathsf B$. This also explains why
possible mixed-degree vectors in the critical kernel cause no
problem for the stated estimate.

For compactly supported $w$, self-adjointness and integration by
parts give
\[
 \int\|(\partial_t+\mathsf B)w\|^2\,\dd t
       =\int\bigl(\|\partial_tw\|^2+\|\mathsf Bw\|^2\bigr)\,\dd t.
\]
Together with \eqref{eq:log-radial-norms} and
\eqref{eq:pure-tangential-coercivity}, this proves the estimate on
the open cone.

Choose a cutoff equal to one far on the cone and zero on the core.
Applying the cone estimate and the first-order Leibniz rule gives
\begin{equation}\label{eq:cap-core-remainder}
 \|v\|\le C\|\rho_Z\mathcal D_Zv\|+C\|v\|_{L^2(K_Z)}
\end{equation}
for a fixed compact set $K_Z$. The completed cap has no $L^2$
harmonic forms in degrees one and two. Indeed its exact-cone-ended
metric is a scattering metric, and
\cite[Theorem~1A]{hausel-hunsicker-mazzeo} identifies its below-middle
$L^2$ cohomology with $H_c^p(T^*S^3;\RR)$. The Thom isomorphism
gives $H_c^p(T^*S^3)=H^{p-3}(S^3)=0$ for $p=1,2$.

If \eqref{eq:weighted-cap} failed, unit compactly supported pure
forms $v_j$ could satisfy $\|\rho_Z\mathcal D_Zv_j\|\to0$.
Interior first-order elliptic estimates and Rellich compactness
give a subsequence converging locally in $L^2$. Its limit is an
$L^2$ harmonic pure form by Fatou and elliptic regularity, hence
zero. The compact term in \eqref{eq:cap-core-remainder} tends to
zero, contradicting the unit norm.
\end{proof}

\subsection{Scaling and preservation of norm}

On the $i$-th scaled cap put
$\rho_{i,\varepsilon}=\varepsilon_i\rho_{Z_i}$; it is the physical
cone radius on the end. Degreewise $L^2$ isometric scaling turns
\eqref{eq:weighted-cap} into
\[
 \|v\|_{L^2(\varepsilon_i^2g_i)}
       \le C\|\rho_{i,\varepsilon}
                    \mathcal D_\varepsilon v\|_{L^2(\varepsilon_i^2g_i)}.
\]
The first-order squared energy scales two powers faster than the mass,
and the squared radius restores those powers.

Choose a radial cutoff $\chi_s$ equal to one on the cap and for
$r\le s$, and zero for $r\ge2s$, with
$|\dd\chi_s|\le C/s$. If $\varepsilon_i$ is sufficiently small
relative to $s$, its support is contained in the part identified with
the completed cap, where
$\rho_{i,\varepsilon}\le Cs$ and
$\rho_{i,\varepsilon}|\dd\chi_s|\le C$ on the shell. Extending
$\chi_su$ by zero and applying the estimate proves
\begin{equation}\label{eq:cap-tail}
 \|u\|^2_{\{i\text{-th cap and }r<s\}}
       \le Cs^2q_\varepsilon(u)
          +C\|u\|^2_{\{s<r<2s\},i},\qquad p=1,2.
\end{equation}
The same constant works for the finitely many fixed caps. Its proof
involves each scale only on its own cap.

\begin{proposition}\label{prop:no-loss-compactness}
Suppose every $\varepsilon_i\to0$, the norms of pure degree-one or
degree-two forms $u_\varepsilon$ are bounded, and
$q_\varepsilon(u_\varepsilon)\to0$. A subsequence converges locally
in $L^2$ to a smooth central $L^2$ closed and coclosed form $u_0$,
and
\begin{equation}\label{eq:no-loss}
 \|u_\varepsilon\|^2-\|u_0\|^2\longrightarrow0.
\end{equation}
The central restrictions extended by zero over missing tips converge
strongly in central $L^2$, and the cap-core norms tend to zero.
For finitely many sequences one can preserve the full Gram matrix
on a common subsequence. For eigenforms with eigenvalues tending
to zero the convergence is smooth on central compact subsets.
\end{proposition}
\begin{proof}
On fixed central compact sets the norm and first-order energy bounds
give local $H^1$ bounds. Rellich compactness and a diagonal choice
give local strong $L^2$ convergence. Vanishing energy gives
$\dd u_0=\delta u_0=0$ distributionally. Elliptic regularity and
Fatou give smoothness and finite $L^2$ norm.

For fixed $s$, apply \eqref{eq:cap-tail} on every end and pass to
the limit:
\[
 \limsup_{\varepsilon\to0}
    \|u_\varepsilon\|^2_{\{\text{all caps and }r<s\}}
       \le C\sum_i\|u_0\|^2_{\{s<r<2s\},i}.
\]
The right side tends to zero as $s\downarrow0$. This and the
local convergence prove strong central convergence, vanishing cap
mass and \eqref{eq:no-loss}. Simultaneous strong convergence and
polarization preserve the Gram matrix of a finite family.
For eigenforms, fixed-exterior elliptic estimates for the Laplacian
upgrade the central convergence to smooth convergence.
\end{proof}

\subsection{The complete model theorem}

We use also the fixed-central Hodge statement from
\cite[Section~5]{wuebben-hodge}: the central $L^2$ closed and coclosed
one-form space is zero, and the corresponding two-form space is the
four-dimensional space of representatives $\eta_D$, $D\in V$.
These representatives lie in the minimal graph domain. This is a
fixed conic assertion; the compactness estimates above supply the
independent-scale step.

\begin{theorem}\label{thm:complete-model}
For fixed $R,A$ there are $L_0,\gamma>0$ and $e(L_*)\to0$ such
that, for every scale vector with $L_*\ge L_0$, the model has the
following properties.

Its harmonic two-form space is marked by all of $K$ and its mass
converges uniformly to $G_R|_K$. There are exactly two positive
two-form eigenvalues below $\gamma$, both coexact, and
\begin{equation}\label{eq:complete-model-rates}
 (1-e(L_*))\nu_j(\varepsilon)
       \le\lambda^+_{2,j}(\varepsilon)
       \le(1+e(L_*))\nu_j(\varepsilon),\qquad j=1,2.
\end{equation}
The first positive one-form and third positive two-form eigenvalues
are bounded below by $\gamma$.

On the full two-dimensional positive small space $S_\varepsilon$,
$\cA_\varepsilon$ is an isomorphism onto $W_R$ and
\begin{equation}\label{eq:model-cluster-mass}
 \left|M_R(\cA_\varepsilon u,\cA_\varepsilon v)
                           -\langle u,v\rangle\right|
       \le e(L_*)\|u\|\|v\|,\qquad
 \|u-U_\varepsilon\cA_\varepsilon u\|
       \le e(L_*)\|u\|.
\end{equation}
There is no ratio restriction among the scales or positive model
eigenvalues.
\end{theorem}
\begin{proof}
If the one-form gap failed, normalized eigenforms with eigenvalues
tending to zero would, by Proposition~\ref{prop:no-loss-compactness},
give a central $L^2$ harmonic one-form of norm one. This contradicts
the fixed-central vanishing. Failure of a uniform threshold would
allow the sequence with $L_*\ge n$ and eigenvalue below $1/n$, so
the contradiction proves a uniform gap over all scale ratios.

The model has $b_2=2$. Proposition~\ref{prop:trial-quotient-mass}
gives two positive coexact eigenvalues bounded by $C/L_*$. If its
third positive two-form eigenvalue could tend to zero, the first
three positive eigenvectors together with two orthonormal harmonic
forms would have vanishing energy. Proposition~\ref{prop:no-loss-compactness}
would preserve five orthonormal vectors in the four-dimensional
central harmonic space, which is impossible. Hence there is a
uniform third-positive-eigenvalue gap. The positive exact two-form
spectrum equals the positive coexact one-form spectrum, with
multiplicities, by the singular values of $\dd$. The one-form gap
therefore makes both small positive two-form spaces coexact.

For $k\in K$, the trial $T_\varepsilon k$ is closed by
\eqref{eq:exact-model-energy} and represents the corresponding
smoothing class. Its harmonic representative $h_\varepsilon(k)$
has norm at most that of the trial:
\[
 \|h_\varepsilon(k)\|^2
       \le\|T_\varepsilon k\|^2=G_R(k,k)+o(1).
\]
Compactness gives a no-loss subsequential central limit $\eta_D$.
Pair with the fixed compactly supported duals
$\Psi_{F,\rho}$, for all $F\in V$ and fixed $\rho$. The class of
the smoothing representative gives
\[
 \int_{N_\varepsilon}h_\varepsilon(k)\wedge\Psi_{F,\rho}
       =G_R(k,F).
\]
Passing to the limit identifies $G_R(D,F)=G_R(k,F)$ for every $F$,
so $D=k$. Norm preservation proves convergence of the harmonic mass.
Finite-dimensional polarization and a contradiction on unit vectors
make it uniform on $K$.

Next take normalized $u_\varepsilon\in S_\varepsilon$ along any
sequence with $L_*\to\infty$. The trial upper bound gives energy
at most $C/L_*$. A no-loss subsequence converges to $\eta_D$ of
norm one. For every $k\in K$, the trial $T_\varepsilon k$ is closed,
so $\langle u_\varepsilon,T_\varepsilon k\rangle=0$. Strong
convergence gives $G_R(D,k)=0$, hence $D\in W_R$. Since the
observation is a bounded fixed operator on the exterior,
$\cA_\varepsilon u_\varepsilon\to D$. Therefore
$M_R(\cA_\varepsilon u_\varepsilon)\to1$.

Every potentially violating sequence has this subsequence, so the
mass convergence is uniform on all unit vectors in $S_\varepsilon$.
Polarization gives the first estimate in
\eqref{eq:model-cluster-mass}; injectivity and dimension make the
observation an isomorphism. The trial
$U_\varepsilon\cA_\varepsilon u_\varepsilon$ has the same strong
central limit and vanishing cap norm. Subtracting and using a
second sequential contradiction gives the reconstruction estimate.

For the span of the first $j$ positive eigenvectors, the weighted
lower bound and mass capture give
\[
 \frac{E_\varepsilon(\cA_\varepsilon u)}
          {M_R(\cA_\varepsilon u)}
       \le
       \frac{1+CL_*^{-1/2}}{1-e_0(L_*)}
       \frac{q_\varepsilon(u)}{\|u\|^2}
       \le(1+o(1))\lambda^+_{2,j}.
\]
Its image has dimension $j$, so generalized min--max gives the
required relative lower bound. The trial upper bound proves the
other side of \eqref{eq:complete-model-rates}. This establishes
the theorem without division by the weakest eigenvalue.
\end{proof}

These model metrics need not be Kähler, and no complex type or
primitivity conclusion has been used. Their constants may depend on
$R,A$. The fixed-central quotient mass converges to $M$ as $R\downarrow0$
by the central harmonic minimization in \cite[Section~5]{wuebben-hodge}.
When the geometric smoothing scales satisfy
$\varepsilon_i=|\tau_i|^{1/3}$ times fixed positive normalization
factors, the relation
\[
 L_i=\frac13|\log|\tau_i||+O_{R,A}(1)
\]
converts each energy coefficient to
$(1+o(1))\,2\pi/|\log|\tau_i||$ multiplicatively. This conversion
alone is not a comparison with an actual Ricci-flat metric.
